\begin{pastexam}[label={px:ncku-114-cce-ds-16}]{成大電通所 114 學年度 單選題第 16 題}
（單選；每題 5 分，答錯倒扣 2 分）

How many recursive calls would the call ``Fibonacci(10)'' generate in the following program?
\begin{minted}[fontsize=\small]{c}
int Fibonacci(int n) {
    if (n < 1) return 0;
    if (n < 3 ) return 1;
\end{minted}
a) 55\quad b) 108\quad c) 110\quad d) 20\quad e) None of the above
\end{pastexam}

\begin{pastsoln}
答案：b) 108。

原卷的程式只印了前三行，少了遞迴那一行和結尾的大括號。 照題意補上最自然的寫法：
\begin{minted}[fontsize=\small]{c}
    return Fibonacci(n - 1) + Fibonacci(n - 2);
}
\end{minted}
設 $C(n)$ 是呼叫 \texttt{Fibonacci(n)} 時，連同這次呼叫本身總共發生的呼叫次數。 $n = 1, 2$ 時直接回傳，$C(1) = C(2) = 1$；
$n \ge 3$ 時呼叫自己一次，再呼叫 $n - 1$ 和 $n - 2$：
\[
C(n) = 1 + C(n - 1) + C(n - 2).
\]
\begin{center}
\begin{tabular}{l*{10}{c}}
\toprule
$n$ & 1 & 2 & 3 & 4 & 5 & 6 & 7 & 8 & 9 & 10 \\
\midrule
$C(n)$ & 1 & 1 & 3 & 5 & 9 & 15 & 25 & 41 & 67 & 109 \\
$F(n)$ & 1 & 1 & 2 & 3 & 5 & 8 & 13 & 21 & 34 & 55 \\
\bottomrule
\end{tabular}
\end{center}
（\texttt{n < 1} 那一行在 $n = 10$ 時用不到：每次遞迴最小只到 1 或 2。） 也可以直接看出 $C(n) = 2F(n) - 1$：令 $D(n) = C(n) + 1$，
則 $D(n) = D(n - 1) + D(n - 2)$、$D(1) = D(2) = 2$，所以 $D(n) = 2F(n)$。

\texttt{Fibonacci(10)} 總共有 $C(10) = 109$ 次呼叫；題目問它「generate」幾次遞迴呼叫，也就是不算最外面那一次，
$109 - 1 = 108$，選 b)。 （a) 55 是函數的回傳值 $F(10)$，不是呼叫次數。） 呼叫次數和 $F(n)$ 同階、呈指數成長，原因是同一個子問題
被重算很多次，這就是資結書第~18.1.1~節說的重疊子問題，用記憶化或由下往上填表可以降到 $\Theta(n)$。
\end{pastsoln}
